\section{Examples along the flow}\label{sec:flow-examples}

We work on a compact curve $X$ of genus at least two. There, stable bundles
of every rank and degree exist. For stable $F_i,F_j$ of equal slope that are
not isomorphic, $h^0(\Hom(F_j,F_i))=0$ and Riemann--Roch gives
$h^1(\Hom(F_j,F_i))=r_ir_j(g(X)-1)>0$. Every graph $\Gamma$ of the kind
considered below is therefore realized by nonzero harmonic entries.

\subsection{All edges tight}

When every edge is tight and $(\Gamma,(r_i))$ is nondegenerate, there is one tight
component and every small eigenvalue decays like $\tau^{-1}$. For two-step
extensions this is the \emph{positive balance condition}. It asks for
positive numbers $\sigma_{ij}$ on the edges with prescribed margins,
\begin{equation}\label{eq:positive-balance}
 \sum_j\sigma_{ij}=\frac{M_iM_J}{2|M|}\quad(i\in I),\qquad
 \sum_i\sigma_{ij}=\frac{M_jM_I}{2|M|}\quad(j\in J),
\end{equation}
where $M_I=\sum_{i\in I}M_i$ and $M_J=\sum_{j\in J}M_j$. Feasibility is a
classical matrix-scaling question \cite{SinkhornKnopp1967,RothblumSchneider1989}.

Condition \eqref{eq:positive-balance} holds exactly when every edge is tight
and $(\Gamma,(r_i))$ is nondegenerate. Let $\bar w$ take the value $-M_J/(2|M|)$ on
$I$ and $M_I/(2|M|)$ on $J$. It is admissible, with equality on every edge,
and \eqref{eq:positive-balance} is the multiplier equation \eqref{eq:kkt}
for $\bar w$ with multipliers $\sigma$. A pair satisfying these optimality
conditions of a convex program characterizes its minimizer, so $\bar w$ is the
weight grading, every edge is tight, and $(\Gamma,(r_i))$ is nondegenerate.
Conversely, if every edge is tight and $\Gamma$ is connected and bipartite, the
differences $w_j-w_i=\frac12$ and $\sum_iM_iw_i=0$ force $w=\bar w$, and
positive multipliers satisfy \eqref{eq:positive-balance}.

\begin{corollary}\label{cor:balanced}
Assume \eqref{eq:positive-balance}. There are unique conductances
$k^*_{ij}=\lVert\beta_{ij}\rVert^2e^{2(y^*_i-y^*_j)}$ with these margins,
normalized by $\sum_iM_iy^*_i=0$. Then
$\tau\lambda_j(h(\tau))\to\rho_j$ for every initial metric, where
$\rho_1\le\cdots\le\rho_{m-1}$ are the positive eigenvalues of the graph
with conductances $k^*$ and masses $M_i$. The spectral projections onto the
eigenvalues with a common limit converge to the
corresponding graph eigenspaces.
\end{corollary}

\begin{proof}
This is Theorem~\ref{thm:complete-spectrum} with $\mathcal T=\Gamma_1$ and $u^*=k^*$.
\end{proof}

\begin{example}\label{ex:four-cycle-flow}
Take four distinct degree-zero line bundles and all four entries
$\beta_{12},\beta_{14},\beta_{32},\beta_{34}$ nonzero, with $V=1$. The
quantity
\[
 \Theta=\frac{\lVert\beta_{12}\rVert^2\lVert\beta_{34}\rVert^2}
             {\lVert\beta_{14}\rVert^2\lVert\beta_{32}\rVert^2}
\]
is invariant under diagonal changes of the factor metrics. The balanced
matrix, with rows $1,3$ and columns $2,4$, is
\[
 \begin{pmatrix}x&y\\y&x\end{pmatrix},\qquad
 x=\frac{\sqrt\Theta}{4(1+\sqrt\Theta)},\qquad
 y=\frac1{4(1+\sqrt\Theta)}.
\]
Its positive graph eigenvalues are $2x$, $2y$ and $\frac12$. Thus every
initial metric gives
\[
 \lim_{\tau\to\infty}\tau(\lambda_1,\lambda_2,\lambda_3)
 =\Bigl(\frac{\min(1,\sqrt\Theta)}{2(1+\sqrt\Theta)},
        \frac{\max(1,\sqrt\Theta)}{2(1+\sqrt\Theta)},\frac12\Bigr).
\]
The constants retain the invariant $\Theta$ of the extension norms, but
not the extension classes themselves. The tight subgraph is a four-cycle,
so $u^*$ depends on the norms. The value $\frac12$ has eigenvector
$w=\frac14(-1,1,-1,1)$.
\end{example}

The three-factor extension $0\to L_1\oplus L_3\to E\to L_2\to0$ with
$\beta_{12},\beta_{32}\ne0$ and $V=1$ is of this kind. Its weight grading
is $(-\frac16,\frac13,-\frac16)$, and both multipliers equal $\frac16$. So
\[
 \lambda_1\sim\frac1{6\tau},\qquad\lambda_2\sim\frac1{2\tau},
\]
with limiting lines $\CC\diag(1,0,-1)$ and $\CC\diag(1,-2,1)$.

\subsection{A leading coefficient depending on the initial metric}

Let $V=1$, and choose pairwise nonisomorphic stable degree-zero bundles
$F_1,F_2,F_3,F_4$ of ranks $(2,1,1,2)$, with nonzero harmonic entries
$\beta_{12},\beta_{32},\beta_{34}$ and no others.

The weight grading is $w=(-\frac16,\frac13,-\frac13,\frac16)$. The edges
$12$ and $34$ are tight with multipliers $\frac13$. The edge $32$ has slack
$\frac13$, hence exponent $\frac43$; the edge $32$ is slack because
$M_1M_4=4>1=M_2M_3$. The multipliers on the tight edges are positive, so
$(\Gamma,(r_i))$ is nondegenerate. This is the region treated in
\cite[(3.24)--(3.26)]{HKKP2023}.

Theorem~\ref{thm:complete-spectrum} gives:
\begin{itemize}
\item two eigenvalues asymptotic to $1/(2\tau)$ (the tight components
 $\{1,2\}$ and $\{3,4\}$, each contributing $\frac12$), with limiting
 subspace $\operatorname{span}\{\diag(1,-2,0,0),\diag(0,0,2,-1)\}$;
\item one eigenvalue $\lambda_1\sim c(h(0))\tau^{-4/3}$, with limiting line
 $\CC\diag(1,1,-1,-1)$ and $c(h(0))=\frac23k^\infty_{32}(h)$.
\end{itemize}
The coefficient is unbounded on the metrics of one bundle. The explicit
family below makes this quantitative.

\begin{proposition}\label{prop:unequal-rank-family}
Normalize the three entries to have norm one, and for $\rho>0$ put
\[
 g_\rho=\diag(\rho^{1/2}\id_{F_1},\rho^{-1/2}\id_{F_2},
              \rho^{1/2}\id_{F_3},\rho^{-1/2}\id_{F_4}),\qquad
 h_\rho=g_\rho^*h_0.
\]
Then the flows starting at $h_\rho$ satisfy
\[
 c(h_\rho)=\frac2{27}\rho^{-2/3}(1+O(\rho)),\qquad\rho\downarrow0.
\]
\end{proposition}

\begin{proof}
Write $k_1=k_{12}$, $k_2=k_{32}$, $k_3=k_{34}$. The masses give
$p=(k_1/2,-k_1-k_2,k_2+k_3,-k_3/2)$ and
\[
 \dot k_1=-3k_1^2-2k_1k_2,\qquad\dot k_2=-2(k_1+k_3)k_2-4k_2^2,\qquad
 \dot k_3=-3k_3^2-2k_2k_3.
\]
The symmetric solution $k_1=k_3=u$, with initial values $u_0$ and
$r_0=k_2(T)/u_0$, is explicit. In the time $s$ with $ds/d\tau=u$ and $s(T)=0$, put
$D(s)=1+2r_0-2r_0e^{-s}$. Then
\[
 r(s)=\frac{r_0e^{-s}}{D(s)},\qquad
 u(s)=\frac{u_0e^{-3s}}{D(s)},\qquad
 \tau-T=\frac1{u_0}\Bigl(\frac{(1+2r_0)(e^{3s}-1)}3-r_0(e^{2s}-1)\Bigr),
\]
so
\[
 \tau^{4/3}k_2\longrightarrow B=\frac{r_0}{3^{4/3}u_0^{1/3}(1+2r_0)^{2/3}}.
\]
For $h_\rho$ all three initial conductances are $\rho^2$, so $u_0=\rho^2$ and
$r_0=1$, and the diagonal flow gives the coefficient
$\frac23B=2/(27\rho^{2/3})$. As in the proof of
Theorem~\ref{thm:complete-spectrum}, compare the flow from $h_\rho$ with the
corrected diagonal metrics along this solution, from time zero. Its initial
discrepancy from $h_\rho$ is $O(\rho^2)$, and its total residual is
$O(\rho)$. The barriers put the actual coefficient within
$e^{\pm C\rho}$ of the diagonal one.
\end{proof}

\subsection{A degenerate example: logarithmic separation}

Let $L_1,\ldots,L_4$ be pairwise nonisomorphic degree-zero line bundles,
with $V=1$, and take the zigzag $0\to L_1\oplus L_3\to E\to L_2\oplus L_4\to0$
with exactly $\beta_{12},\beta_{32},\beta_{34}$ nonzero. The masses are
equal, so $M_1M_4=M_2M_3$ and $(\Gamma,(r_i))$ is degenerate. The middle edge is
tight, but its multiplier is forced to vanish.
Theorem~\ref{thm:complete-spectrum} does not apply. The diagonal flow is
the scalar zigzag of \cite[Example~5.13]{HKKP2023}, and the transfer
theorem still applies.

\begin{theorem}\label{thm:zigzag-flow}
For every smooth initial metric,
\[
 \lambda_1(h(\tau))\sim\frac1{2\tau\log\tau},\qquad
 \lambda_2(h(\tau))\sim\lambda_3(h(\tau))\sim\frac1{2\tau},\qquad
 \liminf_{\tau\to\infty}\lambda_4(h(\tau))>0.
\]
The first line tends to $\CC\diag(1,1,-1,-1)$, and the second and third
together tend to
\[
 \operatorname{span}_{\CC}\{\diag(1,-1,0,0),\diag(0,0,1,-1)\}.
\]
The constants do not depend on the extension forms or the initial
metric.
\end{theorem}

\begin{proof}
With $k_1=k_{12}$, $k_2=k_{32}$, $k_3=k_{34}$, the diagonal flow is
\[
 \dot k_1=-2k_1(2k_1+k_2),\qquad\dot k_2=-2k_2(k_1+2k_2+k_3),\qquad
 \dot k_3=-2k_3(2k_3+k_2).
\]
Put $u=\sqrt{k_1k_3}$, $r=k_2/u$, $z=\sqrt{k_1/k_3}$, $\xi=(z-1)/(z+1)$,
and $ds/d\tau=2u$ with $s(T)=0$. Then
\[
 \frac{dr}{ds}=-r^2,\qquad\frac{d\xi}{ds}=-2\xi,\qquad
 \frac{d\log u}{ds}=-\Bigl(2+\frac{4\xi^2}{1-\xi^2}+r\Bigr).
\]
Hence $r=(s+r_0^{-1})^{-1}$ and $\xi=\xi_0e^{-2s}$ with $|\xi_0|<1$. There
is $C>0$ with $u=Ce^{-2s}(s+r_0^{-1})^{-1}(1+O(e^{-4s}))$, and
$\tau\sim e^{2s}(s+r_0^{-1})/(4C)$. Therefore
\[
 \tau k_1\to\tfrac14,\qquad\tau k_3\to\tfrac14,\qquad
 \tau\log\tau\,k_2\to\tfrac12,
\]
for every positive initial condition. The pairs $\{1,2\}$ and $\{3,4\}$
give two graph eigenvalues asymptotic to $1/(2\tau)$. On vectors constant on
the pairs, the remaining eigenvalue is asymptotic to $k_2$ (quotient masses
$2,2$). Eliminating the internal directions changes it by
$O(k_2^2/\min(k_1,k_3))=o(k_2)$. So
$\nu_1\sim(2\tau\log\tau)^{-1}$ and $\nu_2\sim\nu_3\sim(2\tau)^{-1}$.

Corollary~\ref{cor:coefficient-transfer} applies with
$f=(\tau\log\tau)^{-1}$ and $f=\tau^{-1}$, which proves the eigenvalue
statements. For the first line, a normalized eigensection has Rayleigh
quotient $O((\tau\log\tau)^{-1})$. Along the diagonal gauges $g_\tau$,
Lemma~\ref{lem:instantaneous-edges} gives $k_1,k_3\sim1/(4\tau)$. With the
form bound \eqref{eq:form-lower}, this forces $x_1=x_2$ and $x_3=x_4$ in
every limit. The
complementary two-dimensional space follows by orthogonality. Full-projection
convergence upgrades these identifications to $L^2\to C^k$ convergence.
\end{proof}

The four-factor first eigensection separates $L_1\oplus L_2$ from
$L_3\oplus L_4$, and the faster subspace separates the members of each
pair. Adjoining the identity to the full limiting small subspace recovers
all four summand projections.
